\documentclass[11pt]{article}
\usepackage[margin=1in]{geometry}
\usepackage{amsmath,amssymb,amsthm}
\usepackage{graphicx}
\usepackage[hidelinks]{hyperref}

\newtheorem*{theorem}{Partial Theorem}
\newtheorem*{lemma}{Lemma}

\newcommand{\diam}{\operatorname{diam}}
\newcommand{\dist}{\operatorname{dist}}

\title{Partial Result: Finite Cotype is Necessary for Nonlinear Expander Families}
\author{Automatic Math Discovery Run \texttt{fa\_banach\_001}}
\date{June 18, 2026}

\begin{document}
\maketitle

\section*{Source Problem}

The source paper is Manor Mendel and Assaf Naor,
\emph{Towards a calculus for non-linear spectral gaps}, arXiv:0910.2041
\cite{MN2009}. On PDF page 20, the authors ask for geometric conditions on
metric spaces that admit expander families with respect to those spaces. For
normed spaces they write that it is tempting to conjecture that having such an
expander is equivalent to having finite cotype. In the parsed source this is
\texttt{data/parsed/arxiv\_sources/0910.2041/zigzag-arxiv-v1.tex},
lines 1238--1247.

\begin{center}
\includegraphics[width=0.96\linewidth]{figures/open_problem_crop.png}
\end{center}

\section*{Definitions}

For a finite connected graph \(G=(V,E)\) and a metric space \((Y,d_Y)\), say
that \(G\) has nonlinear Poincare constant at most \(\Gamma\) with respect to
\((Y,d_Y^2)\) if every map \(f:V\to Y\) satisfies
\[
\frac{1}{|V|^2}\sum_{u,v\in V} d_Y(f(u),f(v))^2
\le
\frac{\Gamma}{|E|}\sum_{\{u,v\}\in E}d_Y(f(u),f(v))^2.
\]
A graph family is an expander family with respect to \(Y\) in this sense if
the degrees are uniformly bounded, \(|V|\to\infty\), and the constants
\(\Gamma\) are uniformly bounded.

We use the Maurey--Pisier finite-representability theorem in the following
standard form: a Banach space \(X\) has no finite cotype if and only if, for
every \(n\) and every \(\eta>0\), \(\ell_\infty^n\) embeds linearly into \(X\)
with distortion at most \(1+\eta\).

\section*{Idea of the Proof}

The obstruction is already visible for \(X=\ell_\infty\), as the source notes.
The Maurey--Pisier theorem says that a Banach space with no finite cotype
contains finite-dimensional copies of \(\ell_\infty^n\) at all scales. Every
\(n\)-vertex graph embeds isometrically into \(\ell_\infty^n\) by distance
coordinates. On a bounded-degree graph most vertex pairs are separated by
order \(\log n\), while edges remain length one. The Poincare constant must
therefore grow at least quadratically in \(\log n\), contradicting uniform
boundedness.

\section*{The Necessary Direction}

\begin{theorem}
Let \(X\) be a Banach space. Suppose there are finite connected graphs
\(G_k=(V_k,E_k)\) with uniformly bounded degrees, \(|V_k|\to\infty\), and
\[
\sup_k \gamma(G_k,\|\cdot\|_X^2)<\infty
\]
in the nonlinear Poincare sense above. Then \(X\) has finite cotype.
\end{theorem}

\begin{proof}
We prove the contrapositive. Assume that \(X\) has no finite cotype. Fix
\(\eta\in(0,1)\). By Maurey--Pisier, for every \(n\) there is a linear map
\(T_n:\ell_\infty^n\to X\) such that
\[
\|z\|_\infty\le \|T_n z\|_X\le (1+\eta)\|z\|_\infty
\qquad (z\in\ell_\infty^n).
\]

Let \(G=(V,E)\) be a connected graph on \(n\) vertices with maximum degree at
most \(D\). Define the distance-coordinate embedding
\[
\Phi:V\to \ell_\infty^V,\qquad
\Phi(v)=(d_G(v,w))_{w\in V}.
\]
For \(u,v\in V\), the triangle inequality gives
\(\|\Phi(u)-\Phi(v)\|_\infty\le d_G(u,v)\), while taking \(w=u\) gives the
reverse inequality. Thus \(\Phi\) is an isometric embedding of the graph
metric into \(\ell_\infty^V\). Put \(f=T_n\Phi:V\to X\). Then
\[
d_G(u,v)\le \|f(u)-f(v)\|_X\le (1+\eta)d_G(u,v)
\]
for all \(u,v\in V\). In particular, every edge has \(X\)-length at most
\(1+\eta\).

We also need the elementary bounded-degree distance estimate. There is a
constant \(c_D>0\), depending only on \(D\), such that for every sufficiently
large connected \(n\)-vertex graph of maximum degree at most \(D\),
\[
\frac{1}{n^2}\sum_{u,v\in V}d_G(u,v)^2\ge c_D(\log n)^2.
\]
Indeed, a ball of radius \(r\) in a graph of maximum degree \(D\) has at most
\(C_D(D-1)^r\) vertices when \(D\ge 3\), and at most \(2r+1\) vertices when
\(D\le 2\). Taking \(r\) to be a sufficiently small constant multiple of
\(\log n\), at least half of the vertices are outside \(B(u,r)\) for every
center \(u\), for all large \(n\).

Therefore the Poincare constant of \(G\) with respect to \(X\) satisfies
\[
\gamma(G,\|\cdot\|_X^2)
\ge
\frac{\frac1{n^2}\sum_{u,v}\|f(u)-f(v)\|_X^2}
{\frac1{|E|}\sum_{\{u,v\}\in E}\|f(u)-f(v)\|_X^2}
\ge
\frac{c_D(\log n)^2}{(1+\eta)^2}.
\]
This lower bound tends to infinity with \(n\). Hence no bounded-degree graph
family with \(|V_k|\to\infty\) can have uniformly bounded nonlinear Poincare
constants with respect to \(X\). This proves the contrapositive.
\end{proof}

\section*{Scope and Limitations}

The theorem proves the necessary direction of the source's finite-cotype
equivalence: existence of a normed-space expander family forces finite cotype.
It does not prove the converse, namely that every Banach space of finite cotype
admits such an expander family.

\section*{Novelty and Search Bounds}

The run indexes were searched for arXiv:0910.2041, nonlinear spectral gaps,
finite cotype, \(\ell_\infty\), Maurey--Pisier, and expanders with respect to
Banach spaces. No prior packet for this source question or this necessary
direction was found. The argument is a short standard consequence of
Maurey--Pisier and the distance embedding into \(\ell_\infty^n\), so human
review should treat it as a clarification of a natural half of the conjecture,
not as a solution of the full equivalence.

\section*{Human Review Recommendation}

Check the normalization of the nonlinear Poincare constant against the source's
notation. The proof is invariant under the usual equivalent normalizations, but
the packet states the edge-average form used in the source introduction.

\begin{thebibliography}{9}
\bibitem{MN2009}
M. Mendel and A. Naor,
\emph{Towards a calculus for non-linear spectral gaps}, arXiv:0910.2041, 2009.

\bibitem{MP1976}
B. Maurey and G. Pisier,
\emph{Series de variables aleatoires vectorielles independantes et proprietes
geometriques des espaces de Banach}, Studia Math. 58 (1976), 45--90.
\end{thebibliography}

\end{document}
